% Migration IDs: P-08-001 through P-08-012 and A-08-* in the manifest.
\section{Financial case studies}
\label{sec:financial-case-studies}
\label{sec:financial-compression-audit}

We apply the algorithm suite to two declared tuple-grammar libraries constructed
from licensed-data workflows and exposed through a public-safe aggregate layer.
They are stylized retrospective cases, not observed inventories of a financial
institution. The question is: given the strategy rows, validation profiles,
module tags, and burden index already available at a compression decision, how
much active-maintenance burden can be removed while preserving the declared
source frontier and identity-union closure? The cases do not test a trading
strategy, estimate alpha, compare deployable portfolios, or measure population
resource savings.

\subsection{Protocols and information timing}

The terminal audit makes one retention decision after development and
validation and then evaluates a fixed later period. The annual audit repeats
the decision at annual origins using trailing information through the
preceding year. Their postdecision diagnostics have different definitions and
units, so neither levels nor changes are pooled. Neither design is prospective:
both use endpoint-stable surviving ETF universes drawn from a current licensed
snapshot without revision timestamps. The annual chronology prevents direct
target-year input to a decision but does not remove this population-level
survivorship and snapshot look-ahead.

Both algorithm inputs contain only aggregate strategy identifiers, exact
validation-profile coordinates, the source frontier, module labels, the
identity-closure declaration, and prespecified weights. Dates, prices,
returns, holdings, and held-out outcomes are absent from the compression API.
No held-out quantity can enter weights, preprocessing, applicability,
algorithm choice, warm starts, deletion order, pruning, or the objective.
Only after a selected library passes every exact certificate does a separate
evaluator read the frozen postdecision diagnostic.

For validation return observations \(r_t\) net of the declared turnover cost,
both workflows use the estimated annualized score
\begin{equation}
 \widehat j=\AORFinancialAnnualization\,\overline r-
 \frac{1}{2}(\AORFinancialRiskAversion)(\AORFinancialAnnualization)\,
 \widehat{\operatorname{Var}}(r),
 \label{eq:financial-validation-score}
\end{equation}
where the variance uses the sample denominator when at least two observations
are present. Net returns subtract one-way cost in proportion to turnover. The
primary cost is \AORFinancialPrimaryCostBps\ basis points and the locked
sensitivity grid is \AORFinancialSensitivityBps\
basis points. Thus the frontier coordinates are estimated mean--variance audit
scores in annualized return units under risk-aversion index
\AORFinancialRiskAversion, not true or
preference-free economic opportunities. The original terminal and annual
workflows compare these floating profiles with registered tolerances
\(\AORTerminalFrontierTolerance\) and
\(\AORAnnualFrontierTolerance\); the journal algorithm layer freezes the resulting
profile coordinates in its committed instance artifacts and rechecks those
tables exactly.

The source structures differ materially, as
Table~\ref{tab:aor-financial-instances} shows. In each case, unique tagged-row
carriers force a substantial core. Fixed-point preprocessing then propagates
those selections, deletes satisfied rows and empty contributions, and reduces
the residual cover to an empty problem. Duplicate-coverage and dominance rules
do not drive either reduction. The empty residual is an exact certificate of
the global cover optimum under the human-proved optimum-preserving
preprocessing propositions, followed by exact reconstruction and
original-instance rechecking. No solver is invoked to create that certificate,
and it is not Lean verification.

The module field is a declared tuple grammar: it records signal family,
lookback, rebalance rule, and related construction labels. It does not establish
that every tag corresponds to a separately archived, documented, and reusable
institutional component. Identity closure further treats the union of those
tags as frictionlessly available. Consequently, the case studies validate the
optimization layer for the declared grammar; they do not estimate the cost of
recovering code, data lineage, approvals, or tacit knowledge from a real
organization.

\AORFinancialResourceTable

\subsection{Algorithm comparison}

For each source and weight schedule, the comparison retains the historical
stepwise safe-deletion endpoint and runs heaviest-safe-first, weighted greedy,
weighted greedy followed by reverse deletion, seeded multi-start random
deletion, and the preprocessed MIP workflow. Requirement-mask DP is recorded
as inapplicable because the mandatory-preprocessed requirement count exceeds
its locked limit. A second-solver row is recorded as unavailable because no
second open-source solver is pinned; no proprietary dependency is introduced.
Every feasible endpoint passes exact inactive-retention, tagged-coverage,
frontier, identity-closure, and rational-burden checks.

Table~\ref{tab:aor-financial-algorithm-burdens} shows why safety and burden
optimality must remain separate. The historical stepwise libraries are safe,
but they carry more validation-computation and governance burden than the
empty-residual optima. Heaviest-safe-first attains the certified optimum for
both sources under both displayed nonuniform schedules. Weighted greedy is
close but not uniformly optimal; reverse deletion removes part or all of that
excess. Multi-start deletion does not dominate the historical endpoint in
every source--schedule pair despite every returned endpoint being certified
safe. These outcomes are case-specific and do not overturn the unbounded
worst-case result for heaviest-safe-first.

\AORFinancialAlgorithmTable

\subsection{Locked retrospective weight replay}

The limited robustness replay fixes three defensible schedules. Equal active
weights treat every maintained strategy as one common slot. The
validation-computation schedule uses only the locked signal grammar and
lookback-based computation declared at compression time. The governance
schedule counts documented deviations from a baseline policy template and is
a review-complexity proxy, not an observed monetary cost. All weights are
positive exact integers for active strategies, the inactive strategy has zero
burden, missing grammar fields cause rejection, and held-out outcomes are
prohibited inputs.

This replay is not prospective evidence: the schedules and earlier outcomes
predate the dedicated robustness lock. The lock records a reproducible,
nonadaptive replay subset, algorithm grid, identity-overlap statistic, burden
formulas, tie rule, unique-carrier audit, and reporting schema. It prevents
further outcome-dependent changes within this replay, but it cannot convert a
retrospective analysis into a prospective one. No schedule was searched or
tuned in this replay to favor a reported burden result.

Table~\ref{tab:aor-financial-robustness} gives the exact
global-versus-historical-stepwise comparison. Under equal weights, the
historical endpoints already attain the minimum active count. Under the two
nonuniform schedules, an unchanged retained cardinality can conceal a
material burden gap because different strategies carry different validation
or governance costs. Thus the manager's choice of burden definition affects
which identities are efficient even when the number of retained strategies
does not change.

\AORFinancialRobustnessTable

Identity comparisons support that interpretation. The fixed historical
endpoint is identical across schedules by construction, whereas the adaptive
algorithms can select different tied or near-tied libraries as the schedule
changes. Nevertheless, every feasible endpoint retains every distinct active
unique carrier in its source. Algorithm ranks are stable for the
heaviest-safe-first and exact-preprocessing workflows in these two cases but
vary for several other methods. Complete strategy-identity sets, pairwise
Jaccard fractions, tied ranks, and reconstruction certificates are reported
in Online Resource~1 as public aggregate library labels, not licensed market
rows.

\subsection{Postdecision diagnostics and licensed-data boundary}

Exact closure equality makes the enabled descendant catalog identical across
all safe endpoints. The reported held-out opportunity statistic is the maximum
ex post score over that catalog. Accordingly, its equality across exactly safe
libraries is largely a mechanical consequence of closure equality, not an
independent validation test. The terminal
quantity remains in its locked-period enabled-descendant net-utility unit; the
annual quantity remains a mean next-year enabled-descendant opportunity unit.
They are evaluated after all selections, gaps, ranks, and exclusions are
fixed, and they are never pooled. Because the maximum is evaluated over a
catalog of candidate descendants, it is also exposed to familiar multiple-
testing and selection concerns; we use it only as a mechanism witness and do
not attach inferential or investment-performance meaning to its level.

Figure~\ref{fig:financial-safe-compression} supplies a separate mechanism
check using the historical stepwise endpoint. A frontier-only comparator can
retain a smaller library but loses module closure and lowers the later
enabled-opportunity diagnostic in both audit-specific scales. The safe
endpoint preserves both objects and hence leaves that diagnostic unchanged.
This retrospective pattern illustrates the bridge mechanism within the
declared tuple grammar. Because the diagnostic is an ex post catalog maximum,
the comparison does not show that compression forecasts returns, improves
investment performance, or controls selection bias.

\AORFinancialComparisonFigure

The licensed panels were used locally to reconstruct validation profiles and
to reproduce the committed source frontiers and module unions. Raw CRSP/WRDS
rows and row-level derivatives are neither embedded in the algorithm instances
nor redistributed. Public artifacts contain aggregate strategy labels,
burdens, structural counts, exact certificates, and permitted postdecision
summaries. Researchers require independent licensed access to rerun the
source-data stages.

For research governance, the implication is practical but bounded. A team can
first declare what operating opportunities and generative capabilities must
survive, then attach auditable active-maintenance weights and solve or
approximate the resulting cover. A strategy removed from the maintained
working library may still have to remain in a legal, compliance, or
reproducibility archive. Exact rechecks make the declared preservation promise
independent of the construction method. In these two cases, a historically
reasonable one-at-a-time process leaves avoidable burden under the nonuniform
indices. A new institutional claim requires prospectively recorded inventories,
point-in-time eligibility, documented capability artifacts, and calibrated
maintenance costs rather than more reuse of the same survivor-defined cases.
